\section{Pulsar: Threshold Lattice Signatures}


Pulsar is Lux Network's threshold signature scheme based on Module-LWE
(threshold ML-DSA-65, byte-equal to FIPS-204). It allows
$t$-of-$n$ validators to collaboratively sign messages without any single validator
holding the full secret key.

\subsection{Design Challenges}

Threshold ML-DSA is hard because of rejection sampling: each signer's partial
signature must independently pass the norm bound, and the combined signature must also
pass. Naive approaches have exponentially low success probability as $t$ grows.

Pulsar addresses this through:
\begin{enumerate}[nosep]
  \item \textbf{Expanded masking range}: Each signer uses a larger $\gamma_1$ to ensure
    individual rejection rarely triggers.
  \item \textbf{Aggregation-friendly bounds}: The combined norm bound accounts for the sum
    of $t$ partial signatures.
  \item \textbf{Share refresh}: Proactive secret sharing allows periodic refresh of shares
    without changing the group public key.
\end{enumerate}

\subsection{Key Distribution}

\begin{lstlisting}[language=Go,caption={Pulsar key generation via Shamir's secret sharing over $R_q$}]
func DistributeKey(n, t int) (PublicKey, []SecretShare) {
    // Generate master secret
    s := sampleSecret()

    // Shamir's secret sharing over the ring R_q
    // f(x) = s + a_1*x + ... + a_{t-1}*x^{t-1}
    coeffs := make([]RingElement, t)
    coeffs[0] = s
    for i := 1; i < t; i++ {
        coeffs[i] = sampleRing()
    }

    // Evaluate polynomial at n points
    shares := make([]SecretShare, n)
    for i := 0; i < n; i++ {
        xi := ringFromScalar(uint32(i + 1))
        shares[i] = SecretShare{
            Index: i + 1,
            Value: evaluatePoly(coeffs, xi),
        }
    }

    // Public key from master secret
    pk := computePublicKey(s)
    return pk, shares
}
\end{lstlisting}

%% ============================================================
